% GENERATED FILE. Edit canonical lessons, then run npm run generate:books.
% canonical-source-hash: fnv1a64:b375df380edcd4c1
% canonical-lessons: RU-C155-milyy, RU-C155-umnyy, RU-C155-glupyy, RU-C155-lenivyy, RU-C155-vezhlivyy

\chapter{Nice, Clever, Silly, Lazy, Polite}
\label{ch:ru-nice-clever-silly-lazy-polite}
\hlchaptermodality{155}

\begin{chapteropening}
\textbf{By the end of this chapter:} \emph{I can say nice, clever, silly, lazy and polite.}

Complete the last lesson of chapter 155: I can say nice, clever, silly, lazy and polite.
\end{chapteropening}

\section[\texorpdfstring{\ru{милый} (mílyy)}{милый (mílyy)}]{\ru{милый} (mílyy) — nice, sweet (of a person)}

\label{lesson:RU-C155-milyy}



\begin{quote}
Before the new one: say the Russian for although, then the Russian for otherwise.
\end{quote}

\subsection*{You'll want to know: \ru{милый}}
\textbf{\ru{милый}} (\emph{mílyy}) — \textquotedblleft{}nice, sweet (of a person)\textquotedblright{}.

\begin{grammarlens}[title={What you've built}]
One more word for this chapter.
\end{grammarlens}

\subsection*{Your turn}
\begin{itemize}
\raggedright
  \item \emph{Say it:} \emph{mílyy}
  \item \emph{Say it:} \emph{mílyy}, once more
  \item \emph{Say it:} say \emph{mílyy}
  \item \emph{Recall:} say the Russian for although, then the Russian for otherwise, then say \emph{mílyy} again
\end{itemize}

\begin{tcolorbox}[breakable,colback=teal!4,colframe=teal!35!black,title={Before you move on}]
What does \ru{милый} mean? (\textquotedblleft{}Nice, sweet (of a person)\textquotedblright{}.) Say it once more.
\end{tcolorbox}
\section[\texorpdfstring{\ru{умный} (úmnyy)}{умный (úmnyy)}]{\ru{умный} (úmnyy) — clever}

\label{lesson:RU-C155-umnyy}



\begin{quote}
Before the new one: say the Russian for otherwise, then the Russian for nice.
\end{quote}

\subsection*{You'll want to know: \ru{умный}}
\textbf{\ru{умный}} (\emph{úmnyy}) — \textquotedblleft{}clever\textquotedblright{}.

\begin{grammarlens}[title={What you've built}]
One more word for this chapter.
\end{grammarlens}

\subsection*{Your turn}
\begin{itemize}
\raggedright
  \item \emph{Say it:} \emph{úmnyy}
  \item \emph{Say it:} \emph{úmnyy}, once more
  \item \emph{Say it:} say \emph{úmnyy}
  \item \emph{Recall:} say the Russian for otherwise, then the Russian for nice, then say \emph{úmnyy} again
\end{itemize}

\begin{tcolorbox}[breakable,colback=teal!4,colframe=teal!35!black,title={Before you move on}]
What does \ru{умный} mean? (\textquotedblleft{}Clever\textquotedblright{}.) Say it once more.
\end{tcolorbox}
\section[\texorpdfstring{\ru{глупый} (glúpyy)}{глупый (glúpyy)}]{\ru{глупый} (glúpyy) — silly, stupid}

\label{lesson:RU-C155-glupyy}



\begin{quote}
Before the new one: say the Russian for nice, then the Russian for clever.
\end{quote}

\subsection*{You'll want to know: \ru{глупый}}
\textbf{\ru{глупый}} (\emph{glúpyy}) — \textquotedblleft{}silly, stupid\textquotedblright{}.

\begin{grammarlens}[title={What you've built}]
One more word for this chapter.
\end{grammarlens}

\subsection*{Your turn}
\begin{itemize}
\raggedright
  \item \emph{Say it:} \emph{glúpyy}
  \item \emph{Say it:} \emph{glúpyy}, once more
  \item \emph{Say it:} say \emph{glúpyy}
  \item \emph{Recall:} say the Russian for nice, then the Russian for clever, then say \emph{glúpyy} again
\end{itemize}

\begin{tcolorbox}[breakable,colback=teal!4,colframe=teal!35!black,title={Before you move on}]
What does \ru{глупый} mean? (\textquotedblleft{}Silly, stupid\textquotedblright{}.) Say it once more.
\end{tcolorbox}
\section[\texorpdfstring{\ru{ленивый} (lenívyy)}{ленивый (lenívyy)}]{\ru{ленивый} (lenívyy) — lazy}

\label{lesson:RU-C155-lenivyy}



\begin{quote}
Before the new one: say the Russian for clever, then the Russian for silly.
\end{quote}

\subsection*{You'll want to know: \ru{ленивый}}
\textbf{\ru{ленивый}} (\emph{lenívyy}) — \textquotedblleft{}lazy\textquotedblright{}.

\begin{grammarlens}[title={What you've built}]
One more word for this chapter.
\end{grammarlens}

\subsection*{Your turn}
\begin{itemize}
\raggedright
  \item \emph{Say it:} \emph{lenívyy}
  \item \emph{Say it:} \emph{lenívyy}, once more
  \item \emph{Say it:} say \emph{lenívyy}
  \item \emph{Recall:} say the Russian for clever, then the Russian for silly, then say \emph{lenívyy} again
\end{itemize}

\begin{tcolorbox}[breakable,colback=teal!4,colframe=teal!35!black,title={Before you move on}]
What does \ru{ленивый} mean? (\textquotedblleft{}Lazy\textquotedblright{}.) Say it once more.
\end{tcolorbox}
\section[\texorpdfstring{\ru{вежливый} (vézhlivyy)}{вежливый (vézhlivyy)}]{\ru{вежливый} (vézhlivyy) — polite}

\label{lesson:RU-C155-vezhlivyy}



\begin{quote}
Before the new one: say the Russian for silly, then the Russian for lazy.
\end{quote}

\subsection*{You'll want to know: \ru{вежливый}}
\textbf{\ru{вежливый}} (\emph{vézhlivyy}) — \textquotedblleft{}polite\textquotedblright{}.

\begin{grammarlens}[title={What you've built}]
One more word for this chapter.
\end{grammarlens}

\subsection*{Your turn}
\begin{itemize}
\raggedright
  \item \emph{Say it:} \emph{vézhlivyy}
  \item \emph{Say it:} \emph{vézhlivyy}, once more
  \item \emph{Say it:} say \emph{vézhlivyy}
  \item \emph{Recall:} say the Russian for silly, then the Russian for lazy, then say \emph{vézhlivyy} again
\end{itemize}

\begin{tcolorbox}[breakable,colback=teal!4,colframe=teal!35!black,title={Before you move on}]
What does \ru{вежливый} mean? (\textquotedblleft{}Polite\textquotedblright{}.) Say it once more.
\end{tcolorbox}
